\documentclass{article}

\usepackage[utf8]{inputenc}     % pdfLaTeX
\usepackage[T1]{fontenc}
\usepackage[magyar]{babel}
\usepackage{amsmath}

% Set page size and margins
% Replace `letterpaper' with `a4paper' for UK/EU standard size
\usepackage[a4paper,top=2cm,bottom=2cm,left=3cm,right=3cm,marginparwidth=1.75cm]{geometry}
%for arrows
\usepackage{textcomp}

\title{Mondatszók}
\author{Lengyel Zoltán Péter\thanks{Csukás Ibolya tanítása alapján}}
\date{Legutóbb frissítve: 2026. Március 20.}

\begin{document}

\maketitle
\tableofcontents
\newpage

\begin{center}
    \section{Mondatszók}
\end{center}

\begin{itemize}
    \item Önálló közlésegységek, állandó alakkal bíró egyszavas mondatok
\end{itemize}

\subsection{Indulatszavak}
\begin{itemize}
    \item aj, ó, jé
\end{itemize}

\subsection{Kapcsolattartás mondatszói}
\begin{itemize}
    \item kapcsolatfelvételt, fenntartást megszakítást jelölnek
    \item szervusztok, ühüm, szóval, ugye, pá
\end{itemize}

\subsection{Akaratkifejező mondatszók}
\begin{itemize}
    \item csitt, hoci, sicc
\end{itemize}

\subsection{Rámutató mondatszók}
\begin{itemize}
    \item lám, íme
\end{itemize}

\subsection{Mondatértékű módosítószavak}
\begin{itemize}
    \item eldöntendő kérdésekre válaszolnak
    \item persze, talán
\end{itemize}

\subsection{Hangutánzó mondatszók} 
\begin{itemize}
    \item bumm, miau, hapci
\end{itemize}
\end{document}